a. E Écrivons l'équation de la réaction :

$$
2 \mathrm{CuO}(s)+\mathrm{C}(s) \longrightarrow \mathrm{CO}_2(g)+2 \mathrm{Cu}(s)
$$

- Déterminons les quantités de matière initiales des réactifs :

$$
\begin{aligned}
& n_{\mathrm{CuO}}^i=\frac{m_{\mathrm{CuO}}}{M_{\mathrm{CuO}}}=\frac{1,00}{63,5+16,0} \approx 0,0126 \mathrm{~mol} \\
& n_{\mathrm{C}}^i=\frac{m_{\mathrm{C}}}{M_{\mathrm{C}}}=\frac{0,120}{12}=0,0100 \mathrm{~mol} .
\end{aligned}
$$

- Déterminons les valeurs de $k$ :

$$
\begin{aligned}
& k_{\mathrm{CuO}}=\frac{n_{\mathrm{CuO}}}{2}=\frac{0,0126}{2}=0,0063 \\
& k_{\mathrm{C}}=\frac{n_{\mathrm{C}}}{1}=\frac{0,0100}{1}=0,0100
\end{aligned}
$$

$k_{\mathrm{CuO}}<k_{\mathrm{C}}$, donc $\mathbf{C u O}$ est le réactif limitant.
b. Prenons comme avancement $x(\mathrm{~mol})$ la quantité de matière de dioxyde de carbone formé. Au cours de la transformation, les quantités de matière des réactifs et des produits sont (en mol) :

$$
\begin{gathered}
n_{\mathrm{CuO}}=0,0126-2 x ; \quad n_{\mathrm{C}}=0,0100-x ; \\
n_{\mathrm{CO}_2}=x ; \quad n_{\mathrm{Cu}}=2 x
\end{gathered}
$$


Puisque, dans l'état final, le réactif limitant a disparu, on peut écrire :

$$
\begin{gathered}
n_{\mathrm{CuO}}=0,0126-2 x_{\max }=0 \\
\boldsymbol{x}_{\max }=\frac{0,0126}{2}=\mathbf{0 , 0 0 6 3 ~ m o l}
\end{gathered}
$$

c. $n_{\mathrm{CO}_2}^f=x_{\max }=0,0063 \mathrm{~mol}$
et $n_{\mathrm{Cu}}^f=2 x_{\max }=\mathbf{0 , 0 1 2 6 ~ m o l}$.
